\par\Needspace{6\baselineskip}\medskip
\begin{center}\hypertarget{native-230-article167}{}\pdfbookmark[2]{167.}{article-167}167.\end{center}
Solutionis expositioni signorum praecipuorum indicem praemittimus. Erunt nobis
\par\noindent
\begin{tabular}{@{}ll@{}}
$t,\ t',\ t'',\ t'''$ & quatuor observationum tempora\\
$a,\ a',\ a'',\ a'''$ & corporis coelestis longitudines geocentricae\\
$\beta,\ \beta',\ \beta'',\ \beta'''$ & eiusdem latitudines\\
$r,\ r',\ r'',\ r'''$ & distantiae a Sole\\
$\rho,\ \rho',\ \rho'',\ \rho'''$ & distantiae a terra\\
$l,\ l',\ l'',\ l'''$ & terrae longitudines heliocentricae\\
$B,\ B',\ B'',\ B'''$ & terrae latitudines heliocentricae\\
$R,\ R',\ R'',\ R'''$ & terrae distantiae a Sole
\end{tabular}
\par\noindent\SourceLeaf{231}
$(n01),\ (n12),\ (n23),\ (n02),\ (n13)$ areae duplicatae triangulorum, quae resp.
inter Solem atque corporis coelestis locum primum et secundum,
secundum et tertium, tertium et quartum, primum et tertium, secundum et quartum continentur
\par\noindent
$\dfrac{1}{(\eta01)}\mathpunct{\raisebox{.45ex}{,}}\quad
\dfrac{1}{(\eta12)}\mathpunct{\raisebox{.45ex}{,}}\quad
\dfrac{1}{(\eta23)}$ quotientes e divisione arearum
\newline
$\frac12(n01),\ \frac12(n12),\ \frac12(n23)$
per areas sectorum respondentium oriundi
\[
P'=\frac{(n12)}{(n01)}\mathpunct{\raisebox{.45ex}{,}}\qquad
P''=\frac{(n12)}{(n23)}
\]
\[
\begin{aligned}
Q'&=\left(\frac{(n01)+(n12)}{(n02)}-1\right)r'^3,\\[.5em]
Q''&=\left(\frac{(n12)+(n23)}{(n13)}-1\right)r''^3
\end{aligned}
\]
\noindent $v,\ v',\ v'',\ v'''$ corporis coelestis longitudines in orbita a puncto arbitrario numeratae.
\par
Denique pro observatione secunda et tertia locos heliocentricos terrae in
sphaera coelesti per $A'$, $A''$ denotabimus, locos geocentricos corporis coelestis
per $B'$, $B''$, eiusdemque locos heliocentricos per $C'$, $C''$.
\par
His ita intellectis negotium primum perinde ut in problemate Sect. praec.
(art. 136) consistet in determinatione situs circulorum maximorum $A'C'B'$,
$A''C''B''$, quorum inclinationes ad eclipticam per $\gamma'$, $\gamma''$ designamus: cum hoc
calculo simul iungetur determinatio arcuum $A'B'=\delta'$, $A''B''=\delta''$. Hinc
manifesto erit
\[
\begin{aligned}
r'&=\surd(\rho'\rho'+2\rho'R'\cos\delta'+R'R')\\
r''&=\surd(\rho''\rho''+2\rho''R''\cos\delta''+R''R''),
\end{aligned}
\]
\noindent sive statuendo $\rho'+R'\cos\delta'=x'$, $\rho''+R''\cos\delta''=x''$,
$R'\sin\delta'=a'$, $R''\sin\delta''=a''$:
\[
\begin{aligned}
r'&=\surd(x'x'+a'a')\\
r''&=\surd(x''x''+a''a'').
\end{aligned}
\]
